% !TEX program = pdflatex
\documentclass{article}
\usepackage{amsmath,amssymb,mathtools,bm}
\DeclarePairedDelimiter\abs{\lvert}{\rvert}
\DeclarePairedDelimiter\norm{\lVert}{\rVert}
\DeclareMathOperator*{\argmin}{arg\,min}
\begin{document}
\section{Equations}
\begin{equation}\label{eq:gauss}
\int_{-\infty}^{\infty} e^{-x^2}\,dx = \sqrt{\pi}
\end{equation}
\begin{align}
f(x) &= \sum_{n=0}^{\infty} \frac{f^{(n)}(a)}{n!}(x-a)^n \label{eq:taylor}\\
g(x) &\coloneqq \prod_{i=1}^{N} \left(1 + \frac{x}{i}\right) \nonumber\\
\hat{\bm\theta} &= \argmin_{\bm\theta \in \Theta} \norm*{\mathbf{y} - X\bm\theta}_2^2 + \lambda\abs{\bm\theta}_1
\end{align}
\begin{gather}
a^2 + b^2 = c^2 \\ \nabla \cdot \mathbf{E} = \frac{\rho}{\varepsilon_0}
\end{gather}
\begin{multline}
p(x) = a_0 + a_1x + a_2x^2 + a_3x^3 + a_4x^4 + a_5x^5 \\ + a_6x^6 + a_7x^7 + a_8x^8 + a_9x^9 + a_{10}x^{10}
\end{multline}
\begin{equation}
\begin{split}
(a+b)^3 &= (a+b)(a+b)^2\\ &= a^3 + 3a^2b + 3ab^2 + b^3
\end{split}
\end{equation}
\[
|x| = \begin{dcases} x & x \ge 0,\\ -x & \text{otherwise},\end{dcases}\qquad
A = \begin{pmatrix} a_{11} & \cdots & a_{1n}\\ \vdots & \ddots & \vdots\\ a_{m1} & \cdots & a_{mn}\end{pmatrix},\quad
\det\begin{vmatrix}1&2\\3&4\end{vmatrix} = -2
\]
Symbols: $\mathbb{R}, \mathcal{L}, \mathfrak{g}$, $\mathbb{N}\subseteq\mathbb{Z}\subsetneq\mathbb{Q}$, $\forall\epsilon>0\,\exists\delta$, $\leqslant, \geqslant, \nsim, \therefore, \blacksquare$.
Equations \eqref{eq:gauss} and \eqref{eq:taylor} use $\binom{n}{k}$, $\tfrac12$, $\overbrace{a+b}^{2}$, $\underbrace{c+d}_{3}$, $\xrightarrow[\text{below}]{\text{above}}$.
\begin{align*}
\lim_{n\to\infty}\left(1+\frac{1}{n}\right)^n &= e, & \oint_{\partial\Omega} \omega &= \int_\Omega d\omega.
\end{align*}
\end{document}
